\section{Background}\label{sec:background}

\method uses graph theory to represent both the street network and the sewer network.
In computational terms, the urban mesh can be modeled as a graph $G=(V,E)$, in which the set $V$ represents relevant points in urban space and the set $E$ represents connections between those points.
In the project domain, vertices correspond to intersections, endpoints, curves, intermediate points, collection points, or topographic points of interest.
Edges correspond to street segments or sewer collector segments, with attributes such as length, geometry, road type, and elevation.

This representation enables the use of classical graph properties to interpret the urban mesh.
The degree of a vertex indicates how many connections reach a point and helps identify crossings or bifurcations.
Paths represent possible sequences for flow conduction, while connectivity indicates whether regions of the mesh remain reachable from a discharge or collection point.
In sewer networks, these concepts help preserve coverage, avoid disconnected segments, and identify points where direction changes or maintenance requirements imply physical structures.

The input street network can initially be interpreted as an undirected graph, because streets provide geometric connections between adjacent points without defining the direction of sewage flow.
After topographic analysis, the resulting sewer network is represented as a directed graph.
In this case, each arc indicates a preferred flow direction, ideally from higher elevations to lower elevations.

Gravity is an important physical criterion in sanitary sewer systems.
Therefore, topography directly influences path selection, edge orientation, and the interpretation of high points, low points, and grade breaks.
Since automatically extracted street data may contain redundant intermediate points, the model also requires graph simplification.
The challenge is to reduce redundancy without removing technically relevant elements, such as corners, sharp curves, collection points, and spacing constraints between manholes.
